% Descent Theory Monad-Comonad Correspondence Diagram (Fig. 1.1)
% Canonical TikZ rendering based on canonical.yaml
% Sources: review-02.md, review-03.md, extraction.md
% Diagram elements:
%   C (top-left), M-Mod_C (bottom-center), Desc_C(M) (top-right)
%   (L_M)-Comod_(M-Mod_C) above Desc_C(M), connected by "def." double bar
%   F_M: C -> M-Mod_C, U_M: M-Mod_C -> C
%   F^{L_M}: M-Mod_C -> Desc_C(M), U^{L_M}: Desc_C(M) -> M-Mod_C
%   Q_M: C -> Desc_C(M) (top horizontal)
%   Loop M on C, loop L_M on M-Mod_C
\documentclass[tikz,border=2pt]{standalone}
\usepackage{dzg-tikz}

\begin{document}
% The diagram is a 3x3 matrix with the following layout:
%   row 1:  [empty,          empty,                (L_M)-Comod...]
%   row 2:  [C,              empty,                Desc_C(M)     ]
%   row 3:  [empty,          M-Mod_C,              empty          ]
\begin{tikzcd}[
    row sep=2.5em,
    column sep=5em,
]
  &
  &
  (L_M)\text{-Comod}_{(M\text{-Mod}_{\mathcal{C}})}
    \arrow[d, "\text{def.}", no head, double equal sign distance] \\
  \mathcal{C}
    \arrow[rr, "Q_M"]
    \arrow[dr, "F_M", shift left=0.6ex]
    \arrow[loop above, "M", looseness=8, out=120, in=60]
  &
  &
  \operatorname{Desc}_{\mathcal{C}}(M)
    \arrow[dl, "U^{L_M}", shift left=0.6ex]
  \\
  &
  M\text{-Mod}_{\mathcal{C}}
    \arrow[ul, "U_M", shift left=0.6ex]
    \arrow[ur, "F^{L_M}", shift left=0.6ex]
    \arrow[loop below, "L_M \mathrel{:=} F_MU_M", looseness=8, out=-60, in=-120]
  &
\end{tikzcd}
\end{document}
